\documentclass[12pt]{article}
\usepackage[T1]{fontenc}
\usepackage[utf8]{inputenc}
\usepackage{lmodern}
\usepackage{textcomp}
\usepackage{enumitem}
\usepackage{geometry}
\usepackage{graphicx}
\usepackage{amsmath, amssymb}
\geometry{margin=1in}
\begin{document}
\begin{center}
Economics - Graduate Level Assessment 3 \
Total Marks: 40 \
Answer all questions.
\end{center}

\section*{Multiple Choice (10 marks)}
\begin{enumerate}[label=\arabic*.]
\item The demand curve for a perfectly competitive firm's product is
\begin{enumerate}[label=(\alph*)]
    \item upward sloping and greater than the market demand curve.
    \item kinked at the market price and less than the market demand curve.
    \item perfectly elastic.
    \item downward sloping and equal to the market demand curve.
    \item downward sloping and less than the market demand curve.
\end{enumerate}
\item The Law of Diminishing Marginal Returns is responsible for
\begin{enumerate}[label=(\alph*)]
    \item MC that first falls, but eventually rises, as output increases.
    \item AP that first falls, but eventually rises, as output increases.
    \item TC that first falls, but eventually rises, as output increases.
    \item AVC that first rises, but eventually falls, as output increases.
    \item MP that first falls, but eventually rises, as output increases.
\end{enumerate}
\item Which of the following could be used as a test for autocorrelation up to third order?
\begin{enumerate}[label=(\alph*)]
    \item The Dickey-Fuller test
    \item The Jarque-Bera test
    \item White's test
    \item The Augmented Dickey-Fuller test
    \item The Durbin Watson test
\end{enumerate}
\item According to Say's law
\begin{enumerate}[label=(\alph*)]
    \item when price goes up supply goes up
    \item price and supply have no correlation
    \item when supply goes up, demand goes down
    \item demand creates its own supply
    \item when demand goes up, supply goes down
\end{enumerate}
\item Suppose that there are only two goods: x and y. Which of the following is NOT correct?
\begin{enumerate}[label=(\alph*)]
    \item One can have comparative advantage in producing both goods.
    \item One can have both an absolute advantage and a comparative advantage in producing x.
    \item One can have absolute advantage and no comparative advantage in producing x.
    \item One can have comparative advantage and no absolute advantage in producing x.
\end{enumerate}
\end{enumerate}

\section*{True / False (10 marks)}
\textit{Answer True or False}
\begin{enumerate}[label=\arabic*.]
\setcounter{enumi}{5}
\item True or False: When nominal GDP is $6000 and the GDP deflator is 200, the real GDP is calculated to be $4000.
\item True or False: Total fixed costs (TFC) do not vary with output and are equal to average variable costs (AVC) at all levels of output.
\item True or False: Monopolistically competitive industries are characterized by product differentiation and a lack of perfect substitutability among products.
\item True or False: The appropriate fiscal policy to remedy a recession involves the federal government running a deficit to stimulate economic activity.
\item True or False: An increase in the number of financial intermediaries decreases the diversity of financial products available to consumers.
\end{enumerate}

\section*{Long Form Response (20 marks)}
\textit{Answer each question with a clear and complete explanation.}
\begin{enumerate}[label=\arabic*.]
\setcounter{enumi}{10}
\item Discuss the key factors that contributed to stagflation in the early 1970s, analyzing how
each element interacted within the broader economic landscape.
\item Discuss how the marginal cost equals marginal revenue (MC = MR) rule influences labor
hiring decisions in a perfectly competitive market, emphasizing the relationship between
wage determination and marginal revenue product.
\end{enumerate}

\vfill
\noindent\textit{End of Paper}
\end{document}